\chapter{Methodology: GNNs for Temporal Source Detection}
\label{c:chapter2}
This chapter lays the methodological groundwork for the remainder of the thesis. We begin by formalising the epidemic source detection problem as a probabilistic inference task on temporal contact networks (Section~\ref{st:backtracking}). We then describe how labelled training data are generated through Monte Carlo simulation of the SIR spreading process (Section~\ref{st:data_generation}). Finally, we introduce the three distinct representations of temporal contact information used as inputs to the evaluated model architectures (Section~\ref{st:representation_temp}); the architectures themselves are presented in Chapter~5.

\section{Formalization of the Backtracking Problem}
\label{st:backtracking}
Let $G = (\mathcal{V}, E^t)$ denote a temporal contact network, where $\mathcal{V}$ is a fixed set of $N$ nodes and $E^t = \{(u, v, \tau) : u, v \in \mathcal{V},\, \tau \in \mathbb{N}\}$ is a multiset of timestamped contact events. Each triplet $(u, v, \tau)$ records an undirected interaction between nodes $u$ and $v$ at discrete time step $\tau$. The network can equivalently be viewed as a sequence of snapshot graphs $\mathcal{G} = \{G_0, G_1, \ldots, G_T\}$ over the shared node set $\mathcal{V}$, where $G_\tau = (\mathcal{V}, E_\tau)$ and $E_\tau = \{(u,v) : (u,v,\tau) \in E^t\}$.

An SIR spreading process (defined formally in Chapter~2) is initiated by a single source node $v^* \in \mathcal{V}$ at time $t_0$. The source enters the infectious state $I$ at $t_0$ while all other nodes are susceptible ($S$). The process evolves stochastically along the edges of $\mathcal{G}$ until an observation time $t_{\mathrm{end}} = t_0 + T$, at which point each node $v \in \mathcal{V}$ occupies exactly one of the three compartmental states $X_v \in \{S, I, R\}$. The resulting final snapshot is $S_T = \{C_v : v \in \mathcal{V}\}$, where $C_v \in \{0,1\}^3$ is the one-hot encoding of the state of node $v$ at time $t_{\mathrm{end}}$.

At inference time, two sources of information are available: the final state snapshot $S_T$ and the full temporal contact network $\mathcal{G}$. The exact start time $t_0$ is not known precisely; it is only assumed to lie within a known interval $[t_{\min}, t_{\max}]$. Crucially, the trajectory of the epidemic between $t_0$ and $t_{\mathrm{end}}$ is unobserved — only the terminal state is accessible. This partial observability is the defining characteristic of the backtracking problem and its primary source of difficulty.

The source detection objective is to infer a probability distribution $\hat{p} \in \Delta^{N-1}$ over all nodes, where $\hat{p}_i$ approximates the posterior probability $P(v^* = v_i \mid S_T, \mathcal{G})$. The node $\hat{v}^* = \arg\max_i \hat{p}_i$ is returned as the predicted source. Formally, we seek to learn a parametric function
\[
    f_\theta : (G_a,\, \mathbf{C}) \;\longrightarrow\; \hat{p} \in \Delta^{N-1},
\]
where $G_a$ is the graph structure consumed by the GNN and $\mathbf{C} \in \{0,1\}^{N \times 3}$ is the node feature matrix whose $i$-th row is $C_{v_i}$. The precise form of $G_a$ depends on the temporal representation used and varies across the evaluated models, as discussed in Section~\ref{st:representation_temp}.

The problem is subject to a single-source assumption: exactly one node $v^*$ initiates the outbreak. Under this assumption, a powerful structural constraint holds: any node that is still susceptible at observation time $t_{\mathrm{end}}$ cannot be the source, since the source was infectious at $t_0$ and therefore can never return to the susceptible state. This \textit{expert knowledge mask} eliminates all susceptible nodes from the candidate set and is applied by all GNN architectures evaluated in this thesis (see the \texttt{BacktrackingNetwork} implementation in \texttt{gnn/backtracking\_network.py}). The detection problem is made hard by two compounding factors. First, the SIR spreading process is inherently stochastic: even when the source and the network are held fixed, repeated realisations produce substantially different final snapshots, meaning the mapping from $v^*$ to $S_T$ is many-to-many. Second, given only the final snapshot, many different source nodes may have produced equally plausible outcomes, leading to a multimodal and diffuse posterior $P(v^* \mid S_T, \mathcal{G})$.

\section{Data Generation via Monte Carlo Simulations}
\label{st:data_generation}
Supervised training of a GNN for source detection requires a large collection of labelled instances, each pairing a final SIR snapshot $S_T$ with the corresponding ground-truth source $v^*$. Since labelled epidemiological data of this form are not available at scale in real-world settings, all training data are generated synthetically by simulating the SIR process on the target temporal network.

For each training instance, a source node $v^*$ is selected uniformly at random from those nodes that are active — i.e., present in at least one contact — at a randomly sampled start time $t_0 \in [t_{\min}, t_{\max}]$. The SIR process is then run stochastically on $\mathcal{G}$ from $t_0$ to $t_{\mathrm{end}}$, producing a final state vector $S_T$. The pair $(S_T, v^*)$ constitutes one labelled training instance. Because the spreading process is stochastic, multiple independent realisations from the same source $v^*$ are simulated, exposing the model to the natural variance in outbreak trajectories from a given starting point.

Two distinct sets of simulations are generated per source node. \textit{Ground-truth runs} consist of a large number $M_{\mathrm{gt}}$ of independent realisations per source. These are used exclusively at evaluation time: by averaging predictions over many realisations, the variance in model performance due to stochasticity in $S_T$ is reduced, yielding a stable estimate of source detectability from each possible starting node. \textit{Monte Carlo training runs} consist of a smaller number $M_{\mathrm{mc}}$ of independent realisations per source, drawn from the same distribution. These form the actual training set: a total of $N \times M_{\mathrm{mc}}$ instances, one for each (source, realisation) pair. The separation between ground-truth and training runs ensures that model performance is evaluated on realisations that were never seen during training, providing an unbiased estimate of generalisation. The resulting data are stored as flat binary arrays of type \texttt{int8}, with ground-truth and Monte Carlo states for each compartment ($S$, $I$, $R$) stored in separate files and reshaped to tensors of shape $(N, M, N)$ at load time, as implemented in \texttt{setup/data\_loader.py}.

Not every simulated outbreak provides useful training signal. Outbreaks that fail to spread beyond the immediate neighbourhood of the source, or that extinguish after only a single infection event, offer little information about the network's contact structure and may bias the model towards trivially detectable cases. Simulations are therefore filtered by minimum outbreak size: any realisation in which the number of infected or recovered nodes at $t_{\mathrm{end}}$ falls below a configurable threshold is discarded. This ensures that the training distribution consists of outbreaks that have meaningfully propagated through the network. An additional \textit{possible-sources mask} is computed for each ground-truth run: a node $v$ is flagged as a feasible source candidate for run $r$ if and only if it is non-susceptible in that run, i.e., $C_v \neq S$. This binary mask, of shape $(N, M_{\mathrm{gt}}, N)$, is used to suppress logically impossible source assignments during evaluation.

Data generation is implemented in the \texttt{tsir/} module of this project, which wraps Holme's fast C-based temporal SIR simulation engine~\cite{Holme2020}. The C implementation uses an event-driven algorithm and is substantially faster than equivalent pure-Python realisations, enabling the generation of the large simulation volumes required for supervised training. The C code is invoked from Python via a subprocess interface, and all output is written to disk as binary arrays. Each complete dataset — comprising ground-truth runs, Monte Carlo runs, the temporal network graph, and the possible-sources mask — is logged as a versioned artifact to the Weights and Biases (\textit{W\&B}) experiment tracking platform. Downstream training and evaluation runs reference the artifact by name and version, guaranteeing that every model is trained and evaluated on identical data regardless of when or where the run is executed.

\section{Representation of Temporal Information}
\label{st:representation_temp}
The three GNN architectures evaluated in this thesis operate on the same underlying simulation data but differ fundamentally in how the temporal contact structure of $\mathcal{G}$ is encoded into a graph that can be processed by a GNN. This section describes the three representations in turn. Each is constructed from the temporal graph $\mathcal{G}$ with its associated edge timestamp lists; the corresponding construction code resides in \texttt{gnn/graph\_builder.py}. The choice of representation determines which temporal patterns are preserved, which are discarded, and how the resulting graph interacts with the message-passing operations of the corresponding model architecture.

\subsection{Static Projection (Baseline)}
\label{ssst:staticprojection}
The simplest representation of a temporal contact network is obtained by collapsing the full temporal sequence $\mathcal{G} = \{G_0, G_1, \ldots, G_T\}$ into a single static graph $G_a = (\mathcal{V}, E_a)$, where the aggregated edge set is
\[
    E_a = \bigcup_{\tau=0}^{T} E_\tau.
\]
Every node pair that had at least one contact at any time step appears as a single undirected edge in $G_a$. To retain a measure of contact frequency, edge weights are assigned as $w_{uv} = |\{\tau : (u,v) \in E_\tau\}|$, counting the number of distinct time steps at which the contact was active. When edge weights are used, they are normalised to the unit interval $[0, 1]$ by dividing by the maximum observed count. This representation retains degree structure and contact frequency information but discards all temporal ordering: it is impossible to determine from $G_a$ alone whether a given edge was active early or late in the observation window. The static projection is consumed by the \texttt{StaticGNN} baseline, implemented in \texttt{gnn/static\_gnn.py}, and also provides the structural backbone (without the temporal edge attributes) for the Backtracking Network.

\subsection{Temporal Discretization (Timeslices)}
\label{ssst:temporaldiscret}
The timeslice representation preserves the full temporal ordering of contacts by maintaining a separate edge set for each time step, yielding the sequence $\{E_0, E_1, \ldots, E_T\}$. In practice, many time steps may contain few or no contacts, and processing one GNN layer per original time step is computationally prohibitive for long observation windows. Consecutive time steps are therefore grouped into bins of width $\Delta t$, reducing the effective number of slices to $\lfloor T / \Delta t \rfloor$. All contacts within a window $[\tau \Delta t,\, (\tau+1)\Delta t)$ are merged into a single snapshot edge set $E_\tau^{\Delta t}$. The parameter $\Delta t$ is controlled by the \texttt{group\_by\_time} hyperparameter in the experimental configuration. As $\Delta t \to T$, the timeslice representation converges to the static projection, providing a natural interpolation between fully temporal and fully aggregated representations.

The timeslice sequence $\{E_0^{\Delta t}, \ldots, E_{\lfloor T/\Delta t\rfloor}^{\Delta t}\}$ is consumed by the \texttt{TemporalGNN}, implemented in \texttt{gnn/temporal\_gnn.py}. This architecture applies one dedicated graph convolutional layer (specifically, a GraphSAGE convolution~\cite{Sterchi2025}) per time slice, processing the slices in reverse chronological order — from $t_{\mathrm{end}}$ back towards $t_0$ — in order to propagate information backwards along the direction of potential infection pathways. The reversed traversal is a deliberate design choice motivated by the backtracking objective: features observed at $t_{\mathrm{end}}$ are iteratively propagated to earlier contact neighbourhoods to identify which nodes could plausibly have been active sources.

\subsection{Edge Feature Encoding for Timestamps}
\label{ssst:featureencoding}
The edge texture representation, introduced by Ru et al.~\cite{Ru2023}, encodes the complete temporal activation history of each contact as a fixed-length binary vector attached to the corresponding edge in the aggregated static graph $G_a$. For each edge $(u,v) \in E_a$, the activation vector $X_{(u,v)} \in \{0,1\}^T$ is defined entry-wise as
\[
    \bigl[X_{(u,v)}\bigr]_\tau = \mathbf{1}\bigl[(u,v) \in E_\tau\bigr], \quad \tau = 0, 1, \ldots, T-1.
\]
The full set of activation vectors forms an edge feature matrix $X_a \in \{0,1\}^{|E_a| \times T}$, where $|E_a|$ counts directed edge pairs (both $(u,v)$ and $(v,u)$ are included with the same activation pattern, reflecting the undirected nature of the contacts). This representation preserves the exact temporal pattern of every contact — including the sequential order of activations, the gaps between active periods, and the total contact duration — within a compact fixed-size descriptor that does not depend on the total number of unique time steps visited. The graph $G_a$ together with the matrix $X_a$ constitutes the full input to the \texttt{BacktrackingNetwork}~\cite{Ru2023}, implemented in \texttt{gnn/backtracking\_network.py}. In that architecture, $X_a$ is projected into a hidden embedding space via a learned linear layer before the convolutional message-passing layers, allowing the model to identify temporal patterns in the contact history that are informative for attributing the observed final state $S_T$ to a particular source node.
